\section{Conclusion}
\label{sec: conclusion}

% One paragraph, weighted toward design implications rather than recap: the
% findings are in Section 6 and do not need restating in detail here.  The
% pivot is RQ3's closing claim (the consequential gap is an unreported
% recovery, not an invented action), because that is what earns the specific
% intervention point -- "instrument the divergence" is a conclusion, whereas a
% standalone list of four safeguards reads as future work.
An agent that acts through tools and reports through language leaves two records
of the same episode, and nothing in its architecture holds them together. Across five frontier agents in two environments, the gap between them is universal, uneven by an order of magnitude, unpredictable in direction,
and invisible to factuality. The consequential gaps do not open when the agent
speaks but when execution diverges from what was said---a rejected write, a guessed
identifier, a substituted action---and the narration
continues along the original plan. Oversight therefore belongs at that
divergence rather than at answer time: commitments reconciled against the trace
before the agent summarizes, disclosure of state-changing operations derived
from the trace rather than from narration, and forbidden operations gated by a
monitor that never reads the report. Profiles differ sharply
between agents, so these are a menu rather than a policy, chosen on something
finer than an aggregate rate. Our taxonomy locates these failures
rather than explaining them, and is stable at the gate level but not below it.
An agent's account of its own behavior is a claim like any other, and one that
has, until now, gone unchecked.
